\documentclass[11pt]{amsart}

\usepackage[margin=1.1in]{geometry}
\usepackage{amsmath,amssymb,amsthm,mathtools}
\usepackage{enumitem}
\usepackage{aliascnt}
\usepackage{microtype}
\usepackage{hyperref}
\usepackage[nameinlink,capitalize]{cleveref}

\hypersetup{colorlinks=true,linkcolor=blue,citecolor=blue,urlcolor=blue}
\raggedbottom

\theoremstyle{plain}
\newtheorem{theorem}{Theorem}[section]

\newaliascnt{lemma}{theorem}
\newtheorem{lemma}[lemma]{Lemma}
\aliascntresetthe{lemma}

\newaliascnt{proposition}{theorem}
\newtheorem{proposition}[proposition]{Proposition}
\aliascntresetthe{proposition}

\newaliascnt{corollary}{theorem}
\newtheorem{corollary}[corollary]{Corollary}
\aliascntresetthe{corollary}

\theoremstyle{definition}
\newaliascnt{definition}{theorem}
\newtheorem{definition}[definition]{Definition}
\aliascntresetthe{definition}


\theoremstyle{remark}
\newaliascnt{remark}{theorem}
\newtheorem{remark}[remark]{Remark}
\aliascntresetthe{remark}

\crefname{theorem}{Theorem}{Theorems}
\Crefname{theorem}{Theorem}{Theorems}
\crefname{lemma}{Lemma}{Lemmas}
\Crefname{lemma}{Lemma}{Lemmas}
\crefname{proposition}{Proposition}{Propositions}
\Crefname{proposition}{Proposition}{Propositions}
\crefname{corollary}{Corollary}{Corollaries}
\Crefname{corollary}{Corollary}{Corollaries}
\crefname{definition}{Definition}{Definitions}
\Crefname{definition}{Definition}{Definitions}
\crefname{remark}{Remark}{Remarks}
\Crefname{remark}{Remark}{Remarks}
\title[Local CKN Concentration for Navier--Stokes]
{Local CKN Concentration and the Type I/Type II Dichotomy for Navier--Stokes}
\author{Guillem Cordero}
\date{}

\begin{document}

\begin{abstract}
We prove a local analytic reduction for possible finite-time singularities of
the three-dimensional Navier--Stokes equations.  For a suitable weak solution,
a singular point at time \(T\) must carry a positive amount of
the pressure-normalized Caffarelli--Kohn--Nirenberg critical density, and also
a positive amount of velocity \(L^3\)-concentration, on arbitrarily small
backward parabolic cylinders.  If the critical density vanishes at each point
of the singular set, the solution is regular at time \(T\) by the
epsilon-regularity theorem.  After positive local concentration has been
obtained, the sequence either satisfies a Type I scale-invariant bound or it
constitutes the local Type II alternative.
\end{abstract}

\maketitle

\section{Introduction}

Let \((u,p)\) be a suitable weak solution of the three-dimensional
incompressible Navier--Stokes equations
\[
  \partial_t u+u\cdot\nabla u+\nabla p=\Delta u,
  \qquad \nabla\cdot u=0.
\]
We use the standard normalization in which the viscosity is one.  The weak
solution theory goes back to Leray \cite{Leray1934}.  Suitable weak solutions
and their partial regularity theory were introduced and developed by Scheffer
\cite{Scheffer1976} and by Caffarelli, Kohn, and Nirenberg \cite{CKN1982}; see
also Lin \cite{Lin1998}.  This paper proves the local reduction from a possible
singularity at time \(T\) to the Type I/Type II blowup-rate dichotomy.

Scale-invariant regularity and blow-up criteria form the standard background for
this formulation; see Serrin \cite{Serrin1962}, Giga \cite{Giga1986},
Escauriaza--Seregin--{\v S}verak \cite{ESS2003}, and Seregin
\cite{Seregin2012}.  The proof below uses this literature only for context; its
external analytic input is the Caffarelli--Kohn--Nirenberg epsilon-regularity
theorem.

The argument has four parts.  First, the pressure-normalized critical
quantities \(C\) and \(D\) are finite for suitable weak solutions and invariant
under the usual parabolic scaling.  Second, the Caffarelli--Kohn--Nirenberg
epsilon-regularity theorem implies that vanishing of \(C+D\) at a candidate
singular point gives local regularity, while the velocity-only epsilon
criterion gives a velocity concentration sequence.  Third, a finite-time
singularity therefore produces sequences of shrinking cylinders on which the
critical density, and separately the velocity \(L^3\)-quantity, are bounded
below by positive constants.  Finally, once such a sequence is fixed, it
either satisfies a Type I scale-invariant bound or falls in the local Type II
alternative.

The proof is local in spacetime.  Apart from the standard
Caffarelli--Kohn--Nirenberg epsilon-regularity theorem, all quantities and
alternatives used in the argument are defined below.

\section{Suitable Weak Solutions and Singular Points}

\begin{definition}[Suitable weak solution]
Let \(I\subset\mathbb R\) be an open interval.  A pair \((u,p)\) is a suitable
weak solution on \(\mathbb R^3\times I\) if
\[
  u\in L^\infty_tL^2_{x,\mathrm{loc}}
  \cap L^2_tH^1_{x,\mathrm{loc}},
  \qquad
  p\in L^{3/2}_{\mathrm{loc}},
  \qquad
  \nabla\cdot u=0,
\]
the Navier--Stokes equations hold in the sense of distributions, and the local
energy inequality holds: for almost every \(t_1<t_2\) in \(I\) and every
nonnegative \(\phi\in C_c^\infty(\mathbb R^3\times I)\),
\[
\begin{aligned}
&\int_{\mathbb R^3}\frac{|u(x,t_2)|^2}{2}\phi(x,t_2)\,dx
 +\int_{t_1}^{t_2}\int_{\mathbb R^3}|\nabla u|^2\phi\,dx\,dt \\
&\le
 \int_{\mathbb R^3}\frac{|u(x,t_1)|^2}{2}\phi(x,t_1)\,dx
 +\int_{t_1}^{t_2}\int_{\mathbb R^3}
   \frac{|u|^2}{2}(\partial_t\phi+\Delta\phi)\,dx\,dt \\
&\quad
 +\int_{t_1}^{t_2}\int_{\mathbb R^3}
   \left(\frac{|u|^2}{2}+p\right)u\cdot\nabla\phi\,dx\,dt .
\end{aligned}
\]
The pressure is understood modulo functions of time.
\end{definition}

For \(z_0=(x_0,t_0)\) and \(r>0\), write
\[
  Q_r(z_0)=B_r(x_0)\times(t_0-r^2,t_0).
\]
We also write \(Q_r(x_0,t_0)\) for \(Q_r((x_0,t_0))\).

\begin{definition}[Singular set at a fixed time]
For a time \(T\), define
\[
  \Sigma(T)=
  \{x\in\mathbb R^3:\ u\text{ is not locally bounded in any }
  Q_r(x,T)\}.
\]
Equivalently, \(x\notin\Sigma(T)\) if there are \(r>0\) and \(M<\infty\) such
that \(\|u\|_{L^\infty(Q_r(x,T))}\le M\).
\end{definition}

If a finite-time singularity occurs at time \(T\), then
\(\Sigma(T)\ne\emptyset\) in this local sense.

\section{Parabolic Rescaling and Critical Quantities}

Fix \(z_0=(x_0,T)\).  For \(r>0\), define the parabolic rescaling
\[
  u^{(r)}(y,s)=r\,u(x_0+ry,T+r^2s),
  \qquad
  p^{(r)}(y,s)=r^2p(x_0+ry,T+r^2s).
\]
The rescaled variables are defined on the set of \((y,s)\) for which
\((x_0+ry,T+r^2s)\) belongs to the original spacetime domain.

\begin{proposition}[Invariance under parabolic rescaling]
If \((u,p)\) is a suitable weak solution near \((x_0,T)\), then
\((u^{(r)},p^{(r)})\) is a suitable weak solution on the rescaled domain.  The
local energy inequality is preserved, and the pressure remains defined modulo
functions of the rescaled time.
\end{proposition}

\begin{proof}
The distributional equations and the divergence condition follow from the
change of variables \(x=x_0+ry\), \(t=T+r^2s\).  We spell out the local
energy inequality because this is the only point at which the suitability
condition is not purely formal.

Let \(\psi\in C_c^\infty\) be nonnegative and compactly supported in the
rescaled domain.  Define
\[
  \phi(x,t)=
  \psi\left(\frac{x-x_0}{r},\frac{t-T}{r^2}\right).
\]
If \(y=(x-x_0)/r\) and \(s=(t-T)/r^2\), then
\[
  \partial_t\phi=r^{-2}(\partial_s\psi)(y,s),\qquad
  \nabla_x\phi=r^{-1}(\nabla_y\psi)(y,s),\qquad
  \Delta_x\phi=r^{-2}(\Delta_y\psi)(y,s).
\]
For almost every \(s_1<s_2\), the corresponding times
\(T+r^2s_1<T+r^2s_2\) are admissible in the original local energy inequality.
Apply that inequality for \((u,p)\) to \(\phi\) on
\([T+r^2s_1,T+r^2s_2]\).  After the change of variables
\(x=x_0+ry\), \(t=T+r^2s\), using
\[
  u(x_0+ry,T+r^2s)=r^{-1}u^{(r)}(y,s),\qquad
  p(x_0+ry,T+r^2s)=r^{-2}p^{(r)}(y,s),
\]
all terms contain the same common factor \(r\).  Dividing by this factor gives
exactly
\[
\begin{aligned}
&\int_{\mathbb R^3}\frac{|u^{(r)}(y,s_2)|^2}{2}\psi(y,s_2)\,dy
 +\int_{s_1}^{s_2}\int_{\mathbb R^3}|\nabla_y u^{(r)}|^2\psi\,dy\,ds \\
&\le
 \int_{\mathbb R^3}\frac{|u^{(r)}(y,s_1)|^2}{2}\psi(y,s_1)\,dy
 +\int_{s_1}^{s_2}\int_{\mathbb R^3}
   \frac{|u^{(r)}|^2}{2}(\partial_s\psi+\Delta_y\psi)\,dy\,ds \\
&\quad
 +\int_{s_1}^{s_2}\int_{\mathbb R^3}
   \left(\frac{|u^{(r)}|^2}{2}+p^{(r)}\right)
   u^{(r)}\cdot\nabla_y\psi\,dy\,ds .
\end{aligned}
\]
Thus \((u^{(r)},p^{(r)})\) is suitable on the rescaled domain.  The pressure
transformation is the one dictated by Navier--Stokes scaling, and addition of
a function of \(t\) becomes addition of a function of \(s\).
\end{proof}

\begin{definition}[Pressure-normalized critical quantities]
For \(z_0=(x_0,t_0)\) and \(r>0\), set
\[
  C(z_0,r)=r^{-2}\int_{Q_r(z_0)} |u|^3\,dx\,dt
\]
and
\[
  D(z_0,r)=r^{-2}
  \int_{t_0-r^2}^{t_0}\int_{B_r(x_0)}
  |p(x,t)-(p)_{B_r(x_0)}(t)|^{3/2}\,dx\,dt.
\]
Here \((p)_{B_r(x_0)}(t)\) denotes the spatial mean of \(p(\cdot,t)\) on
\(B_r(x_0)\).
\end{definition}

The subtraction of the spatial mean fixes the pressure ambiguity in the form
used by the Caffarelli--Kohn--Nirenberg criterion.

\begin{proposition}[Local finiteness and scaling]\label{prop:critical-scaling}
Let \((u,p)\) be a suitable weak solution near \(z_0=(x_0,T)\).  Then
\(C(z_0,r)\) and \(D(z_0,r)\) are finite for every sufficiently small \(r\).
Moreover, for every fixed \(R<\infty\),
\[
\begin{aligned}
&\int_{-R^2}^{0}\int_{B_R}
\left(
|u^{(r)}|^3+
|p^{(r)}-(p^{(r)})_{B_R}(s)|^{3/2}
\right)\,dy\,ds \\
&\qquad
=R^2\bigl(C(z_0,rR)+D(z_0,rR)\bigr).
\end{aligned}
\]
\end{proposition}

\begin{proof}
Local finiteness of the velocity term follows from
\[
  u\in L^\infty_tL^2_{x,\mathrm{loc}}
  \cap L^2_tH^1_{x,\mathrm{loc}}
  \subset L^3_{\mathrm{loc}}
\]
on finite cylinders.  Indeed, if \(B\Subset\mathbb R^3\) and
\((a,b)\Subset I\), Sobolev and interpolation give
\[
  \int_a^b \|u(t)\|_{L^3(B)}^3\,dt
  \le
  C
  \left(\operatorname*{ess\,sup}_{a<t<b}\|u(t)\|_{L^2(B)}\right)^{3/2}
  \left(\int_a^b \|u(t)\|_{H^1(B)}^2\,dt\right)^{3/4}
  (b-a)^{1/4}.
\]
Thus \(u\in L^3_{\mathrm{loc}}\).  The pressure term is finite because
\(p\in L^{3/2}_{\mathrm{loc}}\), and subtracting a spatial mean over the same
ball preserves local \(L^{3/2}\) integrability.  More explicitly, for a ball
\(B\),
\[
  |p-(p)_B(t)|^{3/2}\le C\left(|p|^{3/2}+|(p)_B(t)|^{3/2}\right),
\]
while Jensen's inequality gives
\[
  |(p)_B(t)|^{3/2}
  \le \frac1{|B|}\int_B |p(x,t)|^{3/2}\,dx.
\]
After integration over \(B\) and over time, the normalized pressure term is
controlled by the local \(L^{3/2}\) norm of \(p\).

For the scaling identity, use \(x=x_0+ry\), \(t=T+r^2s\).  The velocity term
transforms as
\[
  \int_{-R^2}^{0}\int_{B_R}|u^{(r)}|^3\,dy\,ds
  =r^{-2}\int_{T-(rR)^2}^{T}\int_{B_{rR}(x_0)}|u|^3\,dx\,dt.
\]
The pressure means transform according to
\[
  (p^{(r)})_{B_R}(s)=r^2(p)_{B_{rR}(x_0)}(T+r^2s),
\]
because
\[
  \frac1{|B_R|}\int_{B_R} r^2p(x_0+ry,T+r^2s)\,dy
  =
  \frac{r^2}{|B_{rR}|}\int_{B_{rR}(x_0)}p(x,T+r^2s)\,dx .
\]
Therefore
\[
\begin{aligned}
&\int_{-R^2}^{0}\int_{B_R}
|p^{(r)}-(p^{(r)})_{B_R}(s)|^{3/2}\,dy\,ds \\
&\qquad =
r^{-2}\int_{T-(rR)^2}^{T}\int_{B_{rR}(x_0)}
|p-(p)_{B_{rR}(x_0)}(t)|^{3/2}\,dx\,dt .
\end{aligned}
\]
Adding the velocity and pressure identities gives the displayed formula.
\end{proof}

\section{Epsilon Regularity and the Vanishing Case}

\begin{theorem}[Caffarelli--Kohn--Nirenberg epsilon regularity]
\label{thm:ckn}
There exists a universal constant \(\varepsilon_{\mathrm{CKN}}>0\) with the
following property.  Let \((u,p)\) be a suitable weak solution in \(Q_r(z_0)\).
If
\[
  C(z_0,r)+D(z_0,r)<\varepsilon_{\mathrm{CKN}},
\]
then \(u\) is locally bounded in a smaller cylinder, for instance in
\(Q_{r/2}(z_0)\).  In particular, if \(z_0=(x_0,T)\), then
\(x_0\notin\Sigma(T)\).
\end{theorem}

\begin{proof}
This is the local epsilon-regularity theorem of
Caffarelli--Kohn--Nirenberg \cite{CKN1982}; see also Lin \cite{Lin1998}.  It is
stated here with the pressure normalization used in \(D\).  Equivalent
representatives of the pressure differ by functions of time, and the spatial
mean subtraction on \(B_r(x_0)\) removes this ambiguity.
\end{proof}

\begin{theorem}[Velocity epsilon regularity]\label{thm:velocity-ckn}
There exists a universal constant \(\varepsilon_v>0\) with the following
property.  Let \((u,p)\) be a suitable weak solution in \(Q_r(z_0)\).  If
\[
  C(z_0,r)<\varepsilon_v,
\]
then \(u\) is locally bounded in \(Q_{r/2}(z_0)\).  In particular, if
\(z_0=(x_0,T)\), then \(x_0\notin\Sigma(T)\).
\end{theorem}

\begin{proof}
This is the velocity-only form of the Caffarelli--Kohn--Nirenberg
epsilon-regularity criterion.  It follows from the standard local pressure
decay lemma together with the pressure-normalized criterion
\cref{thm:ckn}; equivalently, it is the \(L^3_{x,t}\)-smallness criterion
proved in the Caffarelli--Kohn--Nirenberg theory and in Lin's direct proof
\cite{CKN1982,Lin1998}.  The statement is scale invariant because \(C\) is
invariant under the parabolic Navier--Stokes scaling.
\end{proof}

\begin{definition}[Vanishing local critical density at time \(T\)]
We say that the local critical density vanishes at time \(T\) if, for every
\(x_0\in\Sigma(T)\),
\[
  \limsup_{r\downarrow0}
  \bigl(C((x_0,T),r)+D((x_0,T),r)\bigr)=0.
\]
This condition is vacuous when \(\Sigma(T)=\emptyset\).
\end{definition}

\begin{proposition}[Equivalent rescaled formulation]
\label{prop:vanishing-rescaled}
The preceding condition is equivalent to the following statement: for every
\(x_0\in\Sigma(T)\) and every fixed \(R<\infty\),
\[
  \lim_{r\downarrow0}
  \int_{-R^2}^{0}\int_{B_R}
  \left(
  |u^{(r)}|^3+
  |p^{(r)}-(p^{(r)})_{B_R}(s)|^{3/2}
  \right)\,dy\,ds=0.
\]
\end{proposition}

\begin{proof}
By \cref{prop:critical-scaling}, the rescaled integral over
\(B_R\times(-R^2,0)\) equals
\[
  R^2\bigl(C((x_0,T),rR)+D((x_0,T),rR)\bigr).
\]
Thus vanishing of \(C+D\) implies vanishing on every fixed rescaled cylinder.
Conversely, taking \(R=1\) gives the original pointwise condition.
\end{proof}

\begin{theorem}[No-concentration regularity]\label{thm:no-concentration}
Let \((u,p)\) be a suitable weak solution on
\(\mathbb R^3\times(T-\delta,T)\) for some \(\delta>0\).  If the local
critical density vanishes at time \(T\), then \(\Sigma(T)=\emptyset\).  Thus no
finite-time singularity occurs at time \(T\) in the no-concentration case.
\end{theorem}

\begin{proof}
Suppose, toward a contradiction, that \(x_0\in\Sigma(T)\).  By the vanishing
hypothesis, choose \(r>0\) such that
\[
  C((x_0,T),r)+D((x_0,T),r)<\varepsilon_{\mathrm{CKN}}.
\]
\Cref{thm:ckn} implies that \(u\) is bounded in \(Q_{r/2}(x_0,T)\), hence
\(x_0\notin\Sigma(T)\).  This contradicts the choice of \(x_0\).  Hence
\(\Sigma(T)=\emptyset\).
\end{proof}

\begin{corollary}[Vanishing rescaled density excludes terminal singularity]
\label{cor:escape}
Let \((u,p)\) be suitable near \((x_0,T)\).  If
\[
  \limsup_{r\downarrow0}
  \bigl(C((x_0,T),r)+D((x_0,T),r)\bigr)=0,
\]
then \(u\) is bounded in some backward cylinder \(Q_\rho(x_0,T)\), and hence
\(x_0\notin\Sigma(T)\).  Equivalently, membership in \(\Sigma(T)\) requires
nonzero pressure-normalized CKN density on a sequence of shrinking backward
parabolic cylinders.
\end{corollary}

\begin{proof}
The hypothesis gives a scale at which \(C+D\) is below the universal threshold
in \cref{thm:ckn}.  The epsilon-regularity theorem gives local boundedness in
a backward cylinder ending at \(T\), so \(x_0\notin\Sigma(T)\).
\end{proof}

\section{Positive Local Concentration}

\begin{lemma}[Failure of vanishing gives a positive concentration sequence]
\label{lem:positive-concentration}
Suppose the local critical density does not vanish at time \(T\).  Then there
exist \(x_0\in\Sigma(T)\), a number \(\eta>0\), and radii \(r_n\downarrow0\)
such that
\[
  C((x_0,T),r_n)+D((x_0,T),r_n)\ge\eta
  \qquad\text{for all }n.
\]
Equivalently, the rescaled sequence around \((x_0,T)\) has nonzero critical
local density on a fixed bounded cylinder.
\end{lemma}

\begin{proof}
Negating the definition of vanishing gives a point \(x_0\in\Sigma(T)\) with
positive limsup of
\[
  C((x_0,T),r)+D((x_0,T),r)
\]
as \(r\downarrow0\).  Choose \(\eta>0\) below that limsup and choose a sequence
\(r_n\downarrow0\) along which the displayed lower bound holds.  The rescaled
formulation follows from \cref{prop:critical-scaling}.
\end{proof}

\begin{definition}[Positive local concentration sequence]
\label{def:positive-concentration-sequence}
Let \(x_0\in\Sigma(T)\).  A sequence \(r_n\downarrow0\) is a positive local
concentration sequence at \((x_0,T)\) if there is \(\eta>0\) such that
\[
  C((x_0,T),r_n)+D((x_0,T),r_n)\ge\eta
  \qquad\text{for all }n.
\]
The corresponding parabolic rescalings are those defined by
\[
  u_n(y,s)=r_nu(x_0+r_ny,T+r_n^2s),
  \qquad
  p_n(y,s)=r_n^2p(x_0+r_ny,T+r_n^2s).
\]
\end{definition}

\begin{proposition}[Persistence of positive concentration]
\label{prop:compactness-persistence}
Let \((u_n,p_n)\) be suitable weak solutions on a fixed cylinder
\(B_R\times(-R^2,0)\).  Suppose that
\[
  \int_{-R^2}^{0}\int_{B_R}
  \left(|u_n|^3+|p_n-(p_n)_{B_R}(s)|^{3/2}\right)\,dy\,ds
  \ge\eta>0
\]
for all \(n\).  Suppose further that, after passing to a subsequence,
\[
  u_n\to U\quad\text{in }L^3(B_R\times(-R^2,0))
\]
and
\[
  p_n-(p_n)_{B_R}(s)\to \Pi
  \quad\text{in }L^{3/2}(B_R\times(-R^2,0)).
\]
Then
\[
  \int_{-R^2}^{0}\int_{B_R}
  \left(|U|^3+|\Pi|^{3/2}\right)\,dy\,ds\ge\eta.
\]
In particular, the limiting pair has nonzero pressure-normalized local
critical density on this cylinder.  If \(\Pi=P-(P)_{B_R}(s)\), this is the
CKN density of \((U,P)\) on \(B_R\times(-R^2,0)\).
\end{proposition}

\begin{proof}
Strong convergence in \(L^3\) gives
\[
  \int_{-R^2}^{0}\int_{B_R}|u_n|^3\,dy\,ds
  \longrightarrow
  \int_{-R^2}^{0}\int_{B_R}|U|^3\,dy\,ds.
\]
Strong convergence in \(L^{3/2}\) gives the analogous convergence of the
normalized pressure integrals.  Passing to the limit in the lower bound yields
the asserted inequality.
\end{proof}

\begin{remark}
\Cref{prop:compactness-persistence} is not used in the proof of the main
dichotomy below.  It is included only as an auxiliary statement for later
compactness arguments, where strong convergence of a rescaled sequence has
already been obtained by other means.
\end{remark}

\begin{proposition}[Local alternative]\label{prop:local-alternative}
Let \((u,p)\) be a suitable weak solution on
\(\mathbb R^3\times(T-\delta,T)\) for some \(\delta>0\).  Exactly one of the
following alternatives holds.
\begin{enumerate}[label=\textup{(\roman*)}]
\item For every \(x_0\in\Sigma(T)\),
\[
  \limsup_{r\downarrow0}
  \bigl(C((x_0,T),r)+D((x_0,T),r)\bigr)=0.
\]
\item There exist \(x_0\in\Sigma(T)\), \(\eta>0\), and \(r_n\downarrow0\) such
that
\[
  C((x_0,T),r_n)+D((x_0,T),r_n)\ge\eta
  \qquad\text{for all }n.
\]
\end{enumerate}
\end{proposition}

\begin{proof}
The first alternative is a universal assertion over the singular set, and the
second is precisely its failure written with a realizing sequence.  If the
first alternative fails, then for some \(x_0\in\Sigma(T)\) the corresponding
limsup is positive; choosing a positive number below this limsup and a
realizing sequence gives the second alternative.  Conversely, the second
alternative contradicts the first.  Hence exactly one alternative holds.
\end{proof}

\begin{theorem}[Singular points carry positive local CKN concentration]
\label{thm:singular-positive}
Let \((u,p)\) be a suitable weak solution on
\(\mathbb R^3\times(T-\delta,T)\).  If \(\Sigma(T)\ne\emptyset\), then there
are \(x_0\in\Sigma(T)\), \(\eta>0\), and \(r_n\downarrow0\) such that
\[
  C((x_0,T),r_n)+D((x_0,T),r_n)\ge\eta
  \qquad\text{for all }n.
\]
\end{theorem}

\begin{proof}
If the first alternative in \cref{prop:local-alternative} held, then
\cref{thm:no-concentration} would imply \(\Sigma(T)=\emptyset\), contrary to
the hypothesis.  Therefore the second alternative in
\cref{prop:local-alternative} holds.
\end{proof}

\begin{corollary}[Singular points carry positive velocity concentration]
\label{cor:velocity-concentration}
Let \((u,p)\) be a suitable weak solution on
\(\mathbb R^3\times(T-\delta,T)\).  If \(x_0\in\Sigma(T)\), then there are
\(\eta_v>0\) and \(r_n\downarrow0\) such that
\[
  C((x_0,T),r_n)\ge\eta_v
  \qquad\text{for all }n.
\]
\end{corollary}

\begin{proof}
If the conclusion failed, then
\[
  \limsup_{r\downarrow0} C((x_0,T),r)=0.
\]
Choose \(r>0\) so small that
\[
  C((x_0,T),r)<\varepsilon_v,
\]
where \(\varepsilon_v\) is the constant in
\cref{thm:velocity-ckn}.  The velocity epsilon-regularity theorem gives
boundedness of \(u\) in \(Q_{r/2}(x_0,T)\), contradicting
\(x_0\in\Sigma(T)\).  Hence the limsup is positive; taking any
\(\eta_v\) below it and passing to a realizing sequence gives the result.
\end{proof}

\section{The Type I/Type II Blowup-Rate Dichotomy}

Let \(x_0\in\Sigma(T)\), \(\eta>0\), and \(r_n\downarrow0\) be given by
\cref{thm:singular-positive}.  Define
\[
  u_n(y,s)=r_nu(x_0+r_ny,T+r_n^2s),
  \qquad
  p_n(y,s)=r_n^2p(x_0+r_ny,T+r_n^2s).
\]
Then \((u_n,p_n)\) are suitable weak solutions on expanding backward
cylinders, and by \cref{prop:critical-scaling} they retain a nonzero
pressure-normalized critical density on a fixed bounded cylinder.

\begin{definition}[Type I concentration sequence]
\label{def:typeI-concentration}
Let \((u,p)\) be a suitable weak solution on
\(\mathbb R^3\times(T-\delta,T)\), and let \(r_n\downarrow0\) be a positive
local concentration sequence at \((x_0,T)\).  The sequence is called Type I
when the concentration point \((x_0,T)\) satisfies a local Type I
scale-invariant bound:
there are \(\rho>0\) and \(M<\infty\), with \(\rho^2<\delta\), such that
\[
  \operatorname*{ess\,sup}_{T-\rho^2<t<T}
  \sqrt{T-t}\,\|u(t)\|_{L^\infty(B_\rho(x_0))}\le M.
\]
The corresponding case is the Type I alternative.
\end{definition}

\begin{definition}[Local Type II alternative]
\label{def:typeII-alternative}
Let \((u,p)\) be a suitable weak solution on
\(\mathbb R^3\times(T-\delta,T)\), and let \(r_n\downarrow0\) be a positive
local concentration sequence at \((x_0,T)\).  The sequence is in the local
Type II alternative when the concentration point \((x_0,T)\) satisfies no
local Type I scale-invariant bound.  Equivalently, for every \(\rho>0\) with
\(\rho^2<\delta\),
\[
  \operatorname*{ess\,sup}_{T-\rho^2<t<T}
  \sqrt{T-t}\,\|u(t)\|_{L^\infty(B_\rho(x_0))}=\infty.
\]
\end{definition}

\begin{theorem}[Local blowup-rate dichotomy]\label{thm:blowup-rate-dichotomy}
Let \((u,p)\) be a suitable weak solution on
\(\mathbb R^3\times(T-\delta,T)\), and suppose \(\Sigma(T)\ne\emptyset\).
Then there exist \(x_0\in\Sigma(T)\),
\(\eta>0\), and \(r_n\downarrow0\) such that
\[
  C((x_0,T),r_n)+D((x_0,T),r_n)\ge\eta
  \qquad\text{for all }n.
\]
The sequence \(r_n\) is a positive local concentration sequence.  It is either a
Type I concentration sequence or a local Type II alternative.  Consequently,
after the no-concentration case has been excluded by epsilon regularity, every
finite-time singular case reaches exactly one of the two blowup-rate
alternatives.
\end{theorem}

\begin{proof}
If the local critical density vanished at every point of \(\Sigma(T)\), then
\cref{thm:no-concentration} would give \(\Sigma(T)=\emptyset\), contrary to the
hypothesis.  Therefore \cref{thm:singular-positive} gives
\(x_0\in\Sigma(T)\), \(\eta>0\), and \(r_n\downarrow0\) with the stated lower
bound.  The rescaled sequence has positive local critical density by
\cref{prop:critical-scaling}.  It either satisfies the Type I bound in
\cref{def:typeI-concentration}, or it does not.  In the latter case it is, by
\cref{def:typeII-alternative}, a local Type II alternative.
\end{proof}

\begin{remark}
The preceding theorem is a local statement.  Its conclusion is the existence of
a positive pressure-normalized CKN concentration sequence and the dichotomy
between a local Type I bound and failure of every such bound at the concentration
point.
\end{remark}

\begin{thebibliography}{9}

\bibitem{CKN1982}
L.~Caffarelli, R.~Kohn, and L.~Nirenberg,
\emph{Partial regularity of suitable weak solutions of the Navier--Stokes
 equations},
Comm. Pure Appl. Math. \textbf{35} (1982), no. 6, 771--831,
doi:\ \nolinkurl{10.1002/cpa.3160350604}.

\bibitem{ESS2003}
L.~Escauriaza, G.~A.~Seregin, and V.~{\v S}verak,
\emph{\(L_{3,\infty}\)-solutions of the Navier--Stokes equations and
 backward uniqueness},
Russian Math. Surveys \textbf{58} (2003), no. 2, 211--250,
doi:\ \nolinkurl{10.1070/RM2003v058n02ABEH000609}.

\bibitem{Giga1986}
Y.~Giga,
\emph{Solutions for semilinear parabolic equations in \(L^p\) and regularity
 of weak solutions of the Navier--Stokes system},
J. Differential Equations \textbf{62} (1986), no. 2, 186--212,
doi:\ \nolinkurl{10.1016/0022-0396(86)90096-3}.

\bibitem{Leray1934}
J.~Leray,
\emph{Sur le mouvement d'un liquide visqueux emplissant l'espace},
Acta Math. \textbf{63} (1934), no. 1, 193--248,
doi:\ \nolinkurl{10.1007/BF02547354}.

\bibitem{Lin1998}
F.-H.~Lin,
\emph{A new proof of the Caffarelli--Kohn--Nirenberg theorem},
Comm. Pure Appl. Math. \textbf{51} (1998), no. 3, 241--257,
doi:\ \nolinkurl{10.1002/(SICI)1097-0312(199803)51:3<241::AID-CPA2>3.0.CO;2-A}.

\bibitem{Scheffer1976}
V.~Scheffer,
\emph{Partial regularity of solutions to the Navier--Stokes equations},
Pacific J. Math. \textbf{66} (1976), no. 2, 535--552,
doi:\ \nolinkurl{10.2140/pjm.1976.66.535}.

\bibitem{Seregin2012}
G.~Seregin,
\emph{A certain necessary condition of potential blow up for Navier--Stokes
 equations},
Comm. Math. Phys. \textbf{312} (2012), no. 3, 833--845,
doi:\ \nolinkurl{10.1007/s00220-011-1391-x}.

\bibitem{Serrin1962}
J.~Serrin,
\emph{On the interior regularity of weak solutions of the Navier--Stokes
 equations},
Arch. Rational Mech. Anal. \textbf{9} (1962), 187--195,
doi:\ \nolinkurl{10.1007/BF00253344}.

\end{thebibliography}

\end{document}
